\documentclass[aspectratio=169]{beamer}
\usetheme{Madrid}
\usepackage[utf8]{inputenc}
\usepackage[T1]{fontenc}
\usepackage{tikz}
\usepackage{amssymb}
\usepackage{graphicx}
\setbeamertemplate{navigation symbols}{}
\setbeamercolor{title}{fg=white,bg=orange!85!red}
\setbeamercolor{frametitle}{fg=white,bg=orange!85!red}
\setbeamercolor{block title}{bg=orange!85!red!20,fg=black}
\setbeamercolor{block body}{bg=orange!85!red!8,fg=black}
\setbeamercolor{item}{fg=orange!85!red}
\setbeamercolor{itemize item}{fg=orange!85!red}
\setbeamertemplate{itemize item}[triangle]

\title{Web Story 4 — One Loop, Many Choices}
\subtitle{A loop that builds page elements — one per item, automatically}
\author{ByteForge --- Grade 8 --- App Builder}
\date{}

\begin{document}

\begin{frame}
    \titlepage
\end{frame}

\begin{frame}{Today's Plan}
    \begin{itemize}
        \item A List of Choices
        \item Numbered Items
        \item Loop Inside a Function
        \item[$\bigstar$] BOSS --- Full Inventory Display
    \end{itemize}
\end{frame}

\begin{frame}{A loop builds elements — it doesn't just repeat text}
    \large
    \begin{itemize}
        \item Each pass through the loop creates a genuinely NEW \texttt{<li>} element and adds it to the page
        \item The array's length directly decides how many list items appear.
    \end{itemize}
\end{frame}


\begin{frame}{The loop counter is real, usable data}
    \large
    \begin{itemize}
        \item The loop's own counter (\texttt{i}) can be used directly in what you build
        \item Numbering items, styling every other row differently, anything depending on POSITION.
    \end{itemize}
\end{frame}


\begin{frame}{A loop wrapped in a function becomes reusable}
    \large
    \begin{itemize}
        \item \texttt{renderList(items)} doesn't care what's inside \texttt{items}
        \item Call it with a shopping list, an enemy list, a chapter list, and the same loop logic handles all of them.
    \end{itemize}
\end{frame}


\begin{frame}{BOSS — array, function, and loop, together}
    \large
    \begin{itemize}
        \item An array holding real data, a function organizing the logic, and a loop turning data into visible elements
        \item Genuinely the core pattern behind almost every real list, menu, or gallery on the web.
    \end{itemize}
\end{frame}

\begin{frame}{Real-World Examples}
    \large
    \begin{itemize}
        \item Every product grid on a shopping site, every social feed, every playlist
        \item All rendered by exactly this "loop through an array, build one element per item" pattern, whether the array has 3 items or 3,000.
    \end{itemize}
\end{frame}


\begin{frame}{Let Them Discover}
    \Large
    Find a real website with a list that clearly scales (a product grid, search results). Do you think the HTML for every single item is hand-written, or generated by a loop like this chapter's? How could you tell?
\end{frame}

\begin{frame}{Your Turn}
    \Large
    Open the student map, find \textbf{Web Story 4 — One Loop, Many Choices}, and work through it
    hands-on --- everything you just saw, they now do themselves.
\end{frame}

\end{document}
